\subsubsection*{1.2.5.}
拓扑斯 $Y\times\mathcal D$ 关于 $Y$ 的函子性质与第 XIII 讲的拓扑斯 $Y\times_S S$ 相同，不再列举。为固定记号，仅记下图
\[
\begin{tikzcd}[row sep=huge, column sep=huge]
Y \arrow[d, ""{name=A}]
  & Y \arrow[d, ""{name=B}] \arrow[r, hook, "i"]
  & Y\times\mathcal D \arrow[d]
  & Y\times\mathcal D^* \arrow[d] \arrow[l, hook', "j"'] \\
(P^t)
  & \{0\} \arrow[r, hook, "i"]
  & \mathcal D
  & \mathcal D^* \arrow[l, hook', "j"']
    \arrow[from=B, to=A, "\mathrm{sp}"']
\end{tikzcd}
\]

\subsubsection*{1.2.6.}
设 $K$ 是 $Y\times\mathcal D$ 上的交换群层复形。如 XIII, 1.4.2 一样，从 $K$ 得到 $K(Y\times\mathcal D^*)$ 中的一个区别三角：
\[
\begin{tikzcd}[row sep=large, column sep=large]
{} & \Phi(K) \arrow[dl] & {} \\
\mathrm{sp}^*K_0 \arrow[rr] & {} & K_\eta \arrow[ul]
\end{tikzcd}
\]
此构造可传至导出范畴，定义一个函子
\[
D(Y\times\mathcal D)
\longrightarrow
(\text{区别三角，所属范畴为 }D(Y\times\mathcal D^*)).
\]

假设 $D$ 有定向（例如是复解析的），并设 $T$ 为 $\pi_1(D^*)$ 的正生成元（1.1.6）。对 $K\in\operatorname{Ob}D(Y\times\mathcal D)$，仿照 XIII, 1.4.3，定义\underline{变差} $V=\operatorname{Var}(T)$，
\[
V:\Phi(K)\longrightarrow K_\eta.
\]
$K_\eta$ 的自同态 $T-\mathrm{Id}$ 是复合
\[
T-\mathrm{Id}:K_\eta\longrightarrow\Phi(K)
\longrightarrow K_\eta.
\]

